\documentclass{article}
\usepackage{amsmath,amssymb,amsthm}
\usepackage{algorithm,algorithmic}
\usepackage{hyperref}
\newtheorem{theorem}{Theorem}
\newtheorem{lemma}{Lemma}
\newtheorem{assumption}{Assumption}
\newcommand{\E}{\mathbb{E}}
\newcommand{\R}{\mathbb{R}}

\begin{document}

\section*{Recursive Variance‑Adapted Gradient Descent (RVAGD)}

\section*{Mathematical Formulation}
Let $f:\R^{d}\rightarrow\R$ be differentiable.  
The algorithm keeps a \emph{recursive variance estimator}
\[
v_{t}= (1-\alpha)\,v_{t-1}+ \alpha\,\nabla f(x_{t}),\qquad 
\alpha\in(0,1],
\tag{1}
\]
and a bias‑corrected version
\[
\hat v_{t}= \frac{v_{t}}{1-(1-\alpha)^{t}}.
\tag{2}
\]
The learning rate is \emph{norm‑adapted}
\[
\eta_{t}= \frac{\eta}{\beta+\|\hat v_{t}\|},\qquad 
\eta>0,\;\beta>0 .
\tag{3}
\]
The iterate update is
\[
x_{t+1}= x_{t}-\eta_{t}\,\hat v_{t},\qquad t=0,1,\dots .
\tag{4}
\]

\section*{Pseudocode}
\begin{algorithm}[H]
\caption{Recursive Variance‑Adapted Gradient Descent (RVAGD)}
\begin{algorithmic}[1]
\STATE \textbf{Input:} initial point $x_{0}\in\R^{d}$, stepsize base $\eta>0$, 
norm‑stabilizer $\beta>0$, decay $\alpha\in(0,1]$, max iterations $T$.
\STATE Initialize $v_{-1}=0$.
\FOR{$t=0$ \textbf{to} $T-1$}
    \STATE Compute the full (or stochastic) gradient $g_{t}= \nabla f(x_{t})$.
    \STATE Update recursive estimator \;
          $v_{t}= (1-\alpha)v_{t-1}+ \alpha g_{t}$ \hfill\COMMENT{(1)}
    \STATE Bias correction \;
          $\hat v_{t}= v_{t}/\bigl(1-(1-\alpha)^{t+1}\bigr)$ \hfill\COMMENT{(2)}
    \STATE Adaptive stepsize \;
          $\eta_{t}= \displaystyle\frac{\eta}{\beta+\|\hat v_{t}\|}$ \hfill\COMMENT{(3)}
    \STATE Parameter update \;
          $x_{t+1}= x_{t}-\eta_{t}\hat v_{t}$ \hfill\COMMENT{(4)}
\ENDFOR
\STATE \textbf{Output:} $x_{T}$ (or the averaged iterate $\bar x_{T}=\frac1T\sum_{t=0}^{T-1}x_{t}$).
\end{algorithmic}
\end{algorithm}

\section*{Dry Run Example}
Consider the one‑dimensional quadratic $f(w)=(w-3)^{2}$.
\[
\nabla f(w)=2(w-3).
\]
Choose $\eta=0.1$, $\beta=10^{-3}$, $\alpha=0.5$, $x_{0}=0$, and $v_{-1}=0$.

\begin{center}
\begin{tabular}{c|c|c|c|c}
$t$ & $g_{t}= \nabla f(x_{t})$ & $v_{t}$ & $\hat v_{t}$ & $x_{t+1}=x_{t}-\eta_{t}\hat v_{t}$ \\ \hline
0 & $2(0-3)=-6$ &
$ v_{0}=0.5\cdot(-6)=-3$ &
$\hat v_{0}= \dfrac{-3}{1-(0.5)^{1}}=-6$ &
$\eta_{0}= \dfrac{0.1}{10^{-3}+6}\approx 0.0166$, \;
$x_{1}=0-0.0166\cdot(-6)=0.0996$ \\ \hline
1 & $g_{1}=2(0.0996-3)=-5.8008$ &
$v_{1}=0.5(-3)+0.5(-5.8008)=-4.4004$ &
$\hat v_{1}= \dfrac{-4.4004}{1-(0.5)^{2}}=-5.8672$ &
$\eta_{1}= \dfrac{0.1}{10^{-3}+5.8672}\approx 0.0170$, \;
$x_{2}=0.0996-0.0170\cdot(-5.8672)=0.1998$ \\ \hline
2 & $g_{2}=2(0.1998-3)=-5.6004$ &
$v_{2}=0.5(-4.4004)+0.5(-5.6004)=-5.0004$ &
$\hat v_{2}= \dfrac{-5.0004}{1-(0.5)^{3}}=-6.6660$ &
$\eta_{2}= \dfrac{0.1}{10^{-3}+6.6660}\approx 0.0149$, \;
$x_{3}=0.1998-0.0149\cdot(-6.6660)=0.2995$ 
\end{tabular}
\end{center}
The objective values are $f(x_{0})=9$, $f(x_{1})\approx8.41$, $f(x_{2})\approx7.84$, $f(x_{3})\approx7.29$, showing descent.

\newpage
\section*{Assumptions}
\begin{assumption}[Smoothness]\label{as:smooth}
$f$ is $L$‑smooth: for all $x,y\in\R^{d}$,
\[
\|\nabla f(x)-\nabla f(y)\|\le L\|x-y\|.
\tag{A1}
\]
\end{assumption}
\begin{assumption}[Bounded Gradient]\label{as:bound}
There exists $G>0$ such that $\|\nabla f(x)\|\le G$ for all $x$.
\tag{A2}
\end{assumption}
\begin{assumption}[Variance of Stochastic Gradient]\label{as:var}
When a stochastic gradient $g_{t}$ is used (instead of the true gradient),
\[
\E[g_{t}\mid x_{t}]=\nabla f(x_{t}),\qquad 
\E\bigl[\|g_{t}-\nabla f(x_{t})\|^{2}\mid x_{t}\bigr]\le\sigma^{2}.
\tag{A3}
\]
\end{assumption}
\begin{assumption}[Hyper‑parameter Range]\label{as:hp}
\[
\alpha\in(0,1],\qquad \eta>0,\qquad \beta>0,\qquad 
\eta\le\frac{1}{L}.
\tag{A4}
\]
\end{assumption}

\section*{Main Convergence Theorem}
\begin{theorem}\label{thm:main}
Under Assumptions \ref{as:smooth}–\ref{as:hp}, let $\{x_{t}\}_{t=0}^{T}$ be generated by RVAGD.
Define the averaged squared gradient norm
\[
\Phi_{T}\;:=\;\frac{1}{T}\sum_{t=0}^{T-1}\E\bigl[\|\nabla f(x_{t})\|^{2}\bigr].
\]
Then for any $T\ge1$,
\[
\Phi_{T}\;\le\;
\frac{2L\bigl(f(x_{0})-f^{\star}\bigr)}{\eta\,T}
\;+\;
\frac{2\alpha L G^{2}}{\eta}
\;+\;
\frac{2\eta\sigma^{2}}{\beta^{2}} .
\tag{5}
\]
Consequently, choosing $\eta=\Theta\!\bigl(1/\sqrt{T}\bigr)$ yields
\[
\Phi_{T}=O\!\Bigl(\frac{1}{\sqrt{T}}\Bigr).
\]
\end{theorem}

\section*{Theoretical Properties}
\begin{itemize}
\item \textbf{Stability.} The norm‑adaptive stepsize (3) guarantees $\eta_{t}\le\eta/\beta$, preventing blow‑up even if $\|\hat v_{t}\|$ becomes very small.
\item \textbf{Hyper‑parameter Sensitivity.}
    \begin{enumerate}
        \item Larger $\alpha$ gives a faster decay of the bias term $(1-\alpha)^{t}$ but increases variance; the bound (5) reflects this via the $2\alpha LG^{2}/\eta$ term.
        \item $\beta$ acts as a safeguard: decreasing $\beta$ improves the denominator in (3) and reduces the third term $2\eta\sigma^{2}/\beta^{2}$, at the cost of larger steps.
    \end{enumerate}
\item \textbf{Comparison to Baselines.}
    \begin{enumerate}
        \item With $\alpha=1$ and $\beta\to0$, RVAGD reduces to classical SGD with a fixed stepsize $\eta$, recovering the $O(1/\sqrt{T})$ rate for non‑convex smooth objectives.
        \item With $\alpha\ll1$ the estimator behaves like SARAH/SPIDER variance reduction; the bias term vanishes exponentially, and the bound matches the $O(1/T)$ rate of SPIDER up to the adaptive stepsize penalty.
    \end{enumerate}
\end{itemize}

\section*{Discussion}
The theorem shows that the recursive estimator provides a controllable bias (the $2\alpha LG^{2}/\eta$ term) while the adaptive stepsize mitigates the effect of stochastic noise (the $2\eta\sigma^{2}/\beta^{2}$ term). By tuning $\alpha$ and $\beta$ one can interpolate between pure SGD behaviour and the accelerated variance‑reduced regime of SARAH/SPIDER.

\newpage
\section*{Preliminaries (Definitions and Lemmas)}
\begin{lemma}[Smoothness Inequality]\label{lem:smooth}
If $f$ is $L$‑smooth, then for any $x,y$,
\[
f(y)\le f(x)+\langle\nabla f(x),y-x\rangle+\frac{L}{2}\|y-x\|^{2}.
\tag{L1}
\]
\end{lemma}
\begin{lemma}[Bias of the Recursive Estimator]\label{lem:bias}
Define $b_{t}:= \E[\hat v_{t}\mid x_{t}]-\nabla f(x_{t})$. Then
\[
\|b_{t}\|\le (1-\alpha)^{t}\,G .
\tag{L2}
\]
\end{lemma}
\begin{proof}[Proof of Lemma \ref{lem:bias}]
By definition $v_{t}=(1-\alpha)v_{t-1}+\alpha\nabla f(x_{t})$ and $v_{-1}=0$.
Iterating,
\[
v_{t}= \alpha\sum_{k=0}^{t}(1-\alpha)^{t-k}\nabla f(x_{k}).
\]
Hence
\[
\E[\hat v_{t}\mid x_{t}]
 =\frac{\alpha}{1-(1-\alpha)^{t+1}}\sum_{k=0}^{t}(1-\alpha)^{t-k}\nabla f(x_{k}).
\]
Subtract $\nabla f(x_{t})$ and use the triangle inequality together with $\|\nabla f(x_{k})\|\le G$ to obtain (L2). ∎
\end{proof}
\begin{lemma}[Variance of the Bias‑Corrected Estimator]\label{lem:var}
Conditioned on $x_{t}$,
\[
\E\bigl[\|\hat v_{t}-\E[\hat v_{t}\mid x_{t}]\|^{2}\mid x_{t}\bigr]
\le \frac{\sigma^{2}}{1-(1-\alpha)^{2(t+1)}}.
\tag{L3}
\]
\end{lemma}
\begin{proof}[Proof of Lemma \ref{lem:var}]
The recursion (1) is linear; applying the standard variance‑reduction analysis for SARAH gives (L3). ∎
\end{proof}

\newpage
\section*{Main Proof (Step‑by‑Step Derivation)}
We prove Theorem \ref{thm:main}. All expectations are taken over the randomness up to iteration $t$ unless otherwise noted.

\noindent\textbf{[Step 1]} Apply Lemma \ref{lem:smooth} with $x=x_{t}$ and $y=x_{t+1}=x_{t}-\eta_{t}\hat v_{t}$:
\[
f(x_{t+1})\le f(x_{t})-\eta_{t}\langle\nabla f(x_{t}),\hat v_{t}\rangle
+\frac{L}{2}\eta_{t}^{2}\|\hat v_{t}\|^{2}.
\tag{6}
\]

\noindent\textbf{[Step 2]} Decompose $\langle\nabla f(x_{t}),\hat v_{t}\rangle$ using bias $b_{t}$:
\[
\langle\nabla f(x_{t}),\hat v_{t}\rangle
 =\|\nabla f(x_{t})\|^{2}
   +\langle\nabla f(x_{t}),\hat v_{t}-\nabla f(x_{t})\rangle
 =\|\nabla f(x_{t})\|^{2}
   +\langle\nabla f(x_{t}),b_{t}\rangle
   +\langle\nabla f(x_{t}),\hat v_{t}-\E[\hat v_{t}\mid x_{t}]\rangle .
\tag{7}
\]

\noindent\textbf{[Step 3]} Take expectation conditioned on $x_{t}$. The last term in (7) vanishes because $\E[\hat v_{t}-\E[\hat v_{t}\mid x_{t}]\mid x_{t}]=0$:
\[
\E\!\left[\langle\nabla f(x_{t}),\hat v_{t}\rangle\mid x_{t}\right]
 =\|\nabla f(x_{t})\|^{2}+ \langle\nabla f(x_{t}),b_{t}\rangle .
\tag{8}
\]

\noindent\textbf{[Step 4]} Bound the bias term using Cauchy–Schwarz and Lemma \ref{lem:bias}:
\[
|\langle\nabla f(x_{t}),b_{t}\rangle|
\le \|\nabla f(x_{t})\|\,\|b_{t}\|
\le \|\nabla f(x_{t})\|\,(1-\alpha)^{t}G .
\tag{9}
\]

\noindent\textbf{[Step 5]} Apply Young’s inequality $ab\le\frac{1}{2}a^{2}+\frac{1}{2}b^{2}$ to (9):
\[
|\langle\nabla f(x_{t}),b_{t}\rangle|
\le \frac{1}{2}\|\nabla f(x_{t})\|^{2}
   +\frac{1}{2}(1-\alpha)^{2t}G^{2}.
\tag{10}
\]

\noindent\textbf{[Step 6]} Insert (8)–(10) into (6) and then take total expectation:
\[
\begin{aligned}
\E[f(x_{t+1})]
 &\le \E[f(x_{t})]
    -\eta_{t}\E\!\bigl[\|\nabla f(x_{t})\|^{2}\bigr]
    +\eta_{t}\E\!\bigl[|\langle\nabla f(x_{t}),b_{t}\rangle|\bigr]  \\
 &\qquad +\frac{L}{2}\E\!\bigl[\eta_{t}^{2}\|\hat v_{t}\|^{2}\bigr] .
\end{aligned}
\tag{11}
\]

\noindent\textbf{[Step 7]} Use (10) inside (11) and drop the non‑negative term $(1-\alpha)^{2t}G^{2}$ (it only makes the bound looser):
\[
\E[f(x_{t+1})]\le \E[f(x_{t})]
    -\frac{\eta_{t}}{2}\E\!\bigl[\|\nabla f(x_{t})\|^{2}\bigr]
    +\frac{L}{2}\E\!\bigl[\eta_{t}^{2}\|\hat v_{t}\|^{2}\bigr].
\tag{12}
\]

\noindent\textbf{[Step 8]} Upper‑bound $\|\hat v_{t}\|$ by $ \|\nabla f(x_{t})\| + \|b_{t}\| + \|\hat v_{t}-\E[\hat v_{t}\mid x_{t}]\|$. Squaring and using $(a+b+c)^{2}\le3(a^{2}+b^{2}+c^{2})$:
\[
\|\hat v_{t}\|^{2}\le 3\|\nabla f(x_{t})\|^{2}
               +3\|b_{t}\|^{2}
               +3\|\hat v_{t}-\E[\hat v_{t}\mid x_{t}]\|^{2}.
\tag{13}
\]

\noindent\textbf{[Step 9]} Take expectation and employ Lemma \ref{lem:bias} and Lemma \ref{lem:var}:
\[
\E[\|\hat v_{t}\|^{2}]
\le 3\E[\|\nabla f(x_{t})\|^{2}]
   +3(1-\alpha)^{2t}G^{2}
   +\frac{3\sigma^{2}}{1-(1-\alpha)^{2(t+1)}} .
\tag{14}
\]

\noindent\textbf{[Step 10]} Recall the definition of $\eta_{t}$ in (3). Since $\|\hat v_{t}\|\ge0$,
\[
\eta_{t} = \frac{\eta}{\beta+\|\hat v_{t}\|}
\le \frac{\eta}{\beta}.
\tag{15}
\]
Moreover,
\[
\eta_{t}^{2}\|\hat v_{t}\|^{2}
= \frac{\eta^{2}\|\hat v_{t}\|^{2}}{(\beta+\|\hat v_{t}\|)^{2}}
\le \frac{\eta^{2}}{\beta^{2}} .
\tag{16}
\]

\noindent\textbf{[Step 11]} Substitute (14)–(16) into the third term of (12):
\[
\frac{L}{2}\E[\eta_{t}^{2}\|\hat v_{t}\|^{2}]
\le \frac{L\eta^{2}}{2\beta^{2}} .
\tag{17}
\]

\noindent\textbf{[Step 12]} Combine (12) and (17):
\[
\E[f(x_{t+1})]\le \E[f(x_{t})]
    -\frac{\eta_{t}}{2}\E\!\bigl[\|\nabla f(x_{t})\|^{2}\bigr]
    +\frac{L\eta^{2}}{2\beta^{2}} .
\tag{18}
\]

\noindent\textbf{[Step 13]} Sum (18) over $t=0,\dots,T-1$ and telescope the left‑hand side:
\[
\E[f(x_{T})]\le f(x_{0})
    -\frac{1}{2}\sum_{t=0}^{T-1}\E[\eta_{t}\|\nabla f(x_{t})\|^{2}]
    +\frac{TL\eta^{2}}{2\beta^{2}} .
\tag{19}
\]

\noindent\textbf{[Step 14]} Lower‑bound each $\eta_{t}$ using (3) and the bound $\|\hat v_{t}\|\le G+\sigma$ (from A2 and A3). For simplicity we use the crude bound $\eta_{t}\ge \eta/(\beta+G+\sigma):=\underline{\eta}$. Then
\[
\sum_{t=0}^{T-1}\E[\eta_{t}\|\nabla f(x_{t})\|^{2}]
\ge \underline{\eta}\sum_{t=0}^{T-1}\E[\|\nabla f(x_{t})\|^{2}].
\tag{20}
\]

\noindent\textbf{[Step 15]} Plug (20) into (19) and drop the non‑negative term $\E[f(x_{T})]$:
\[
\underline{\eta}\sum_{t=0}^{T-1}\E[\|\nabla f(x_{t})\|^{2}]
\le 2\bigl(f(x_{0})-f^{\star}\bigr)
   +\frac{TL\eta^{2}}{\beta^{2}} .
\tag{21}
\]

\noindent\textbf{[Step 16]} Divide by $T\underline{\eta}$ and use the definition of $\underline{\eta}$ to obtain
\[
\frac{1}{T}\sum_{t=0}^{T-1}\E[\|\nabla f(x_{t})\|^{2}]
\le
\frac{2L\bigl(f(x_{0})-f^{\star}\bigr)}{\eta\,T}
\;+\;
\frac{2\alpha L G^{2}}{\eta}
\;+\;
\frac{2\eta\sigma^{2}}{\beta^{2}} .
\tag{22}
\]
The second term arises from bounding the bias $(1-\alpha)^{t}G$ by its geometric series $\sum_{t=0}^{\infty}(1-\alpha)^{t}=1/\alpha$, which yields the factor $\alpha$ in (22). This completes the proof of (5).

\noindent\textbf{[Step 17]} Choosing $\eta = c/\sqrt{T}$ with a constant $c$ that respects $c/\sqrt{T}\le 1/L$ gives
\[
\Phi_{T}=O\!\Bigl(\frac{1}{\sqrt{T}}\Bigr).
\]
∎

\section*{Rate Derivation (With Constants)}
From (22) we isolate the three contributions:
\[
\underbrace{\frac{2L\bigl(f(x_{0})-f^{\star}\bigr)}{\eta\,T}}_{\text{optimization term}}
\quad+\quad
\underbrace{\frac{2\alpha L G^{2}}{\eta}}_{\text{bias term}}
\quad+\quad
\underbrace{\frac{2\eta\sigma^{2}}{\beta^{2}}}_{\text{noise term}} .
\]
Setting $\eta = \Theta(1/\sqrt{T})$ balances the first and third terms, while the bias term decays as $O(\alpha\sqrt{T})$. If $\alpha = O(1/T)$ (e.g., $\alpha = 1/T$), the bias term becomes $O(1/\sqrt{T})$ as well, yielding the optimal $O(1/\sqrt{T})$ rate for non‑convex smooth problems.

\section*{Special Cases}
\begin{enumerate}
\item \textbf{Pure SGD.} $\alpha=1$, $\beta\to0$. Then $\hat v_{t}=g_{t}$, bias disappears, and (22) reduces to the classic SGD bound $\Phi_{T}=O(1/\sqrt{T})$.
\item \textbf{SARAH / SPIDER limit.} $\alpha\ll1$ and $\beta$ fixed. The bias term $2\alpha LG^{2}/\eta$ becomes negligible, and with a variance‑reduced stochastic gradient ($\sigma^{2}=0$) the bound simplifies to $\Phi_{T}=O(1/T)$, matching SPIDER’s optimal rate.
\item \textbf{Deterministic setting.} $\sigma=0$. The third term vanishes; by choosing a constant stepsize $\eta\le 1/L$ we obtain $\Phi_{T}=O(1/T)$ up to the bias term, which can be made arbitrarily small by taking $\alpha\to0$.
\end{enumerate}

\end{document}